\input{BeamerSettings}

\definecolor{codegreen}{rgb}{0,0.6,0}
\definecolor{codegray}{rgb}{0.5,0.5,0.5}
\definecolor{codepurple}{rgb}{0.58,0,0.82}
\definecolor{backcolour}{rgb}{0.95,0.95,0.92}

\lstdefinestyle{mystyle}{
    backgroundcolor=\color{backcolour},
    commentstyle=\color{codegreen},
    keywordstyle=\color{magenta},
    numberstyle=\tiny\color{codegray},
    stringstyle=\color{codepurple},
    basicstyle=\ttfamily\tiny,
    breakatwhitespace=false,
    breaklines=true,
    captionpos=b,
    keepspaces=true,
    numbers=left,
    numbersep=5pt,
    showspaces=false,
    showstringspaces=false,
    showtabs=false,
    tabsize=2,
    xleftmargin=0.05\textwidth,
    xrightmargin=0.45\textwidth
}

% Style for full-width code blocks (no side margins)
\lstdefinestyle{fullwidth}{
    backgroundcolor=\color{backcolour},
    commentstyle=\color{codegreen},
    keywordstyle=\color{magenta},
    numberstyle=\tiny\color{codegray},
    stringstyle=\color{codepurple},
    basicstyle=\ttfamily\tiny,
    breakatwhitespace=false,
    breaklines=true,
    captionpos=b,
    keepspaces=true,
    numbers=left,
    numbersep=5pt,
    showspaces=false,
    showstringspaces=false,
    showtabs=false,
    tabsize=2,
    xleftmargin=0pt,
    xrightmargin=0pt
}

\AtBeginSection[]{
  \begin{frame}[noframenumbering]
  \vfill
  \centering
  \begin{block}{}
    \usebeamerfont{title}\insertsectionhead\par
  \end{block}
  \vfill
  \end{frame}
}

\newcommand{\inlinecode}[1]{\colorbox{backcolour}{\lstinline{#1}}}

\lstset{style=mystyle}

\author[RMR]{\textbf{Ryan~M.~Richard}}
\title[Software Testing]{Software Testing and PyTest}
\institute[Ames]{
	Scientist and Adjunct Assistant Professor of Chemistry\\
	Ames National Laboratory and Iowa State University\\
	Ames, IA, USA
}

\date[6-8-2026]{2026 SIMCODES Bootcamp}

\begin{document}
\TitleSlide


\section{Terminology}

\TwoColumnSlide{Terminology}
{\BulletList{
	\item<1-> \textbf{Bug}. When software behaves unexpectedly.
	\item<1->Term origin:
	\begin{itemize}
		\item<1-> Apocryphal: Grace Hopper found a moth preventing her program from running.
		\item<1-> Truth: It's actually older, e.g.,  Thomas Edison used it to refer to his inventions failing.
	\end{itemize}
	\item<2> \textbf{Debugging}. Act of finding bugs.
}}
{
\only<1>{\AddPicture{moth.png}}
\only<2>{\AddPicture{debug.jpg}}
}
{}

\TwoColumnSlide{Terminology}
{\BulletList{
	\item \textbf{Test Case}. The \textit{environment}, inputs, outputs, needed to verify correctness.
	\begin{itemize}
		\item Usually check one thing, and only one thing.
		\item Want to check two things? Use two test cases!
	\end{itemize}
	\item \textbf{Test Suite}. A collection of test cases.
	\item \textbf{Testing Framework}. Software for automatically managing test suites.
}}
{\AddPicture{test_suite.png}}
{}

\TwoColumnSlide{Terminology}
{\BulletList{
	\item \textbf{Regression}. When a change to working software breaks it.
	\item \textbf{Regression Testing}. Running a test suite to make sure changes do
	 not break the software.
}}
{\AddPicture{regression.jpg}}
{}

\TwoColumnSlide{Types of Tests Cases}
{\BulletList{
\item \textbf{Unit tests}: tests a single function or class in isolation.
\item \textbf{Integration tests}: tests that multiple components work together.
\item \textbf{End-to-end tests}: tests the full system from a user's perspective.
\item This tutorial focuses on \textbf{unit testing}.
}}
{\AddPicture{unit_test.png}}
{}

%%%%%%%%%%%%%%%%%%%%%%%%%%%%%%%%%%%%%%%%%%%%%%%%%%%
\section{Getting Started with PyTest}
%%%%%%%%%%%%%%%%%%%%%%%%%%%%%%%%%%%%%%%%%%%%%%%%%%%

\TwoColumnSlide{Motivation}
{\BulletList{
	\item Software bugs can be costly or even dangerous.
	\item Manual testing does not scale to large projects.
	\item Automated tests let you:
	\begin{itemize}
		\item Catch bugs early and cheaply.
		\item Detect regressions when you modify code.
		\item Document expected behavior for collaborators.
		\item Develop with confidence.
	\end{itemize}
	\item Rule of thumb: the earlier a bug is found, the easier to fix.
}}
{\AddPicture{software_fail.jpg}}
{}

\TwoColumnSlide{What is PyTest?}
{\BulletList{
	\item Usually also provide features to manage the test suite.
	\item Many Python testing frameworks.
	\item IMHO PyTest is the easiest.
	\item Will get far just using \texttt{assert} statements.
}}
{\AddPicture{pytest.jpg}}
{}

\begin{frame}[fragile,squeeze]{Installing PyTest}
\setlength{\leftmargini}{0.5cm}
\begin{columns}[c,t,totalwidth=1.0\textwidth]
\begin{column}[c]{0.49\linewidth}
\begin{block}{}
\begin{itemize}
\item Install once per Python environment.
\item \textbf{Recommended:} install inside your virtual or conda environment.
\item You can also run PyTest as a module: \texttt{python -m pytest}
\end{itemize}
\end{block}
\end{column}
\begin{column}[c]{0.49\linewidth}
Install with \texttt{pip} (standard Python):
\begin{lstlisting}[style=fullwidth,language=bash,numbers=none]
pip install pytest
\end{lstlisting}
Install with \texttt{conda} (Anaconda/miniconda):
\begin{lstlisting}[style=fullwidth,language=bash,numbers=none]
conda install pytest
\end{lstlisting}
Verify the installation:
\begin{lstlisting}[style=fullwidth,language=bash,numbers=none]
pytest --version
# pytest 8.x.y
\end{lstlisting}
\end{column}
\end{columns}
\end{frame}

\begin{frame}[fragile,squeeze]{Test Discovery: PyTest's Naming Rules}
\setlength{\leftmargini}{0.5cm}
\begin{columns}[c,t,totalwidth=1.0\textwidth]
\begin{column}[c]{0.49\linewidth}
\begin{block}{}
\begin{itemize}
\item PyTest \textbf{automatically finds} tests --- no configuration needed.
\item It searches the current directory (and subdirectories) for:
\begin{itemize}
\item Files named \texttt{test\_*.py} or \texttt{*\_test.py}.
\item Functions named \texttt{test\_*()} inside those files.
\item Methods named \texttt{test\_*()} inside \texttt{Test*} classes.
\end{itemize}
\item Follow these conventions and PyTest finds everything automatically.
\end{itemize}
\end{block}
\end{column}
\begin{column}[c]{0.49\linewidth}
\begin{lstlisting}[style=fullwidth,language=python,numbers=none]
# Discovered: file is named test_math.py

# Discovered: starts with test_
def test_add():
    ...

# NOT discovered: missing test_ prefix
def check_add():
    ...

# NOT discovered: too generic
def test():
    ...
\end{lstlisting}
\end{column}
\end{columns}
\end{frame}

\begin{frame}[fragile,squeeze]{Running Tests}
\setlength{\leftmargini}{0.5cm}
\begin{columns}[c,t,totalwidth=1.0\textwidth]
\begin{column}[c]{0.49\linewidth}
\begin{block}{}
\begin{itemize}
\item Run all tests in the current directory tree.
\item Verbose output (show each test name).
\item Run tests in a specific file.
\item Run a single test by name.
\item Run tests matching a keyword.
\item Stop after the first failure.
\end{itemize}
\end{block}
\end{column}
\begin{column}[c]{0.49\linewidth}
\begin{lstlisting}[style=fullwidth,language=bash,numbers=none]
pytest

pytest -v

pytest test_math_utils.py

pytest test_math_utils.py::test_add

pytest -k "add"

pytest -x
\end{lstlisting}
\end{column}
\end{columns}
\end{frame}

%%%%%%%%%%%%%%%%%%%%%%%%%%%%%%%%%%%%%%%%%%%%%%%%%%%
\section{Writing Your First Test}
%%%%%%%%%%%%%%%%%%%%%%%%%%%%%%%%%%%%%%%%%%%%%%%%%%%

\begin{frame}[fragile,squeeze]{Example Module: \texttt{math\_utils.py}}
\begin{block}{}
Create a file \texttt{math\_utils.py}. We will use these three functions as running examples.
\end{block}
\begin{lstlisting}[language=python,numbers=left,xleftmargin=0.05\textwidth,xrightmargin=0.05\textwidth]
import math

def add(a, b):
    """Return the sum of a and b."""
    return a + b

def divide(a, b):
    """Return a divided by b. Raises ValueError if b is zero."""
    if b == 0:
        raise ValueError("Cannot divide by zero")
    return a / b

def circle_area(radius):
    """Return the area of a circle with the given radius."""
    if radius < 0:
        raise ValueError("Radius cannot be negative")
    return math.pi * radius ** 2
\end{lstlisting}
\end{frame}

\begin{frame}[fragile,squeeze]{Your First Test File: \texttt{test\_math\_utils.py}}
\setlength{\leftmargini}{0.5cm}
\begin{columns}[c,t,totalwidth=1.0\textwidth]
\begin{column}[c]{0.49\linewidth}
\begin{lstlisting}[style=fullwidth,language=python,numbers=left]
# test_math_utils.py
from math_utils import add

def test_add_positive_numbers():
    assert add(2, 3) == 5

def test_add_negative_numbers():
    assert add(-1, -1) == -2

def test_add_with_zero():
    assert add(0, 5) == 5
\end{lstlisting}
\end{column}
\begin{column}[c]{0.49\linewidth}
\begin{block}{}
\begin{itemize}
\item Import the function you want to test.
\item Each test function must start with \texttt{test\_}.
\item Use Python's built-in \texttt{assert} statement.
\item One logical check per test is best practice.
\item Descriptive names act as documentation.
\end{itemize}
\end{block}
Run the tests:
\begin{lstlisting}[style=fullwidth,language=bash,numbers=none]
pytest -v test_math_utils.py
\end{lstlisting}
\end{column}
\end{columns}
\end{frame}

\begin{frame}[fragile,squeeze]{Reading PyTest Output}
\setlength{\leftmargini}{0.5cm}
\begin{columns}[c,t,totalwidth=1.0\textwidth]
\begin{column}[c]{0.49\linewidth}
\begin{block}{All tests pass}
\begin{lstlisting}[style=fullwidth,language=bash,numbers=none]
test_math_utils.py::test_add_positive PASSED
test_math_utils.py::test_add_negative PASSED
test_math_utils.py::test_add_with_zero PASSED

3 passed in 0.01s
\end{lstlisting}
\end{block}
\end{column}
\begin{column}[c]{0.49\linewidth}
\begin{block}{A failing test}
\begin{lstlisting}[style=fullwidth,language=bash,numbers=none]
FAILED test_math_utils.py::test_add_positive
AssertionError: assert 4 == 5
  where 4 = add(2, 3)

1 failed in 0.02s
\end{lstlisting}
\end{block}
\begin{block}{}
\begin{itemize}
\item \texttt{.} = passed \quad \texttt{F} = failed
\item \texttt{E} = error \quad \texttt{s} = skipped
\item PyTest shows exactly what went wrong and where.
\end{itemize}
\end{block}
\end{column}
\end{columns}
\end{frame}

\TwoColumnSlide{Hands-On: Write Your First Test}
{\BulletList{
\item Create \texttt{math\_utils.py} with the three functions from the example.
\item Create \texttt{test\_math\_utils.py}.
\item Write at least one test for \texttt{add()}.
\item Run \inlinecode{pytest -v} and confirm all tests pass.
\item Deliberately break a test (e.g., \inlinecode{assert add(2,3) == 99}) and observe the output.
\item Fix it and re-run.
}}
{\AddPicture{write_unit_test.png}}
{}

\TwoColumnSlide{Hands-On: Exploring PyTest Options}
{\BulletList{
\item Experiment with these flags:
\begin{itemize}
\item \inlinecode{pytest -v} --- verbose mode
\item \inlinecode{pytest -x} --- stop on first failure
\item \inlinecode{pytest -k "zero"} --- run tests matching keyword
\end{itemize}
\item Goal: get comfortable with the edit-test-fix cycle.
}}
{\AddPicture{coding_cycle.jpg}}
{}

\begin{frame}[fragile,squeeze]{The Floating Point Problem}
\setlength{\leftmargini}{0.25cm}
\begin{columns}[c,t,totalwidth=1.0\textwidth]
\begin{column}[c]{0.45\linewidth}
\begin{block}{}
\begin{itemize}
\item Floating point arithmetic is not exact.
\item Never use \texttt{==} to compare two floats in tests.
\item Use \texttt{pytest.approx} instead.
\item Default tolerance: $10^{-6}$ (relative).
\item Can set absolute or relative tolerance.
\end{itemize}
\end{block}
\end{column}
\begin{column}[c]{0.48\linewidth}
\begin{lstlisting}[style=fullwidth,language=python,numbers=left]
import pytest, math
from math_utils import circle_area

# BAD: may fail due to floating point
def test_area_bad():
    assert circle_area(1.0) == math.pi

# GOOD: use pytest.approx
def test_area():
    assert circle_area(1.0) == pytest.approx(math.pi)

# Custom absolute tolerance
def test_area_loose():
    assert circle_area(1.0) == pytest.approx(math.pi,
                                              abs=1e-4)
\end{lstlisting}
\end{column}
\end{columns}
\end{frame}


\begin{frame}[fragile,squeeze]{Testing That Exceptions Are Raised}
\setlength{\leftmargini}{0.5cm}
\begin{columns}[c,t,totalwidth=1.0\textwidth]
\begin{column}[c]{0.49\linewidth}
\begin{lstlisting}[style=fullwidth,language=python,numbers=left]
import pytest
from math_utils import divide, circle_area

def test_divide_by_zero():
    with pytest.raises(ValueError):
        divide(10, 0)

def test_divide_error_message():
    with pytest.raises(ValueError,
                       match="Cannot divide by zero"):
        divide(10, 0)

def test_negative_radius():
    with pytest.raises(ValueError):
        circle_area(-1.0)

def test_divide_does_not_raise():
    # No exception should be raised for valid input
    result = divide(10, 2)
    assert result == pytest.approx(5.0)
\end{lstlisting}
\end{column}
\begin{column}[c]{0.49\linewidth}
\begin{block}{}
\begin{itemize}
\item Use \texttt{pytest.raises()}.
\item The test \textbf{passes} if the expected exception is raised.
\item The test \textbf{fails} if the exception is NOT raised.
\item The \texttt{match} argument checks the error message (regex).
\item Test that your code fails gracefully, not just that it succeeds!
\end{itemize}
\end{block}
\end{column}
\end{columns}
\end{frame}

%%%%%%%%%%%%%%%%%%%%%%%%%%%%%%%%%%%%%%%%%%%%%%%%%%%
\section{Fixtures}
%%%%%%%%%%%%%%%%%%%%%%%%%%%%%%%%%%%%%%%%%%%%%%%%%%%

\TwoColumnSlide{What Are Fixtures?}
{\BulletList{
	\item A \textbf{fixture} is reusable setup code shared across tests.
	\item Common uses:
	\begin{itemize}
		\item Saves time/code creating objects that multiple tests need.
		\item Opening and closing files.
		\item Establishing a known state before each test.
		\item Setup/teardown that needs to occur.
	\end{itemize}
	\item Avoids copy-paste of setup code.
	\item Reduces the code needed to write tests.
}}
{\AddPicture{principles.jpg}}
{}

\begin{frame}[fragile,squeeze]{Using Fixtures}
\setlength{\leftmargini}{0.5cm}
\begin{columns}[c,t,totalwidth=1.0\textwidth]
\begin{column}[c]{0.49\linewidth}
\begin{lstlisting}[style=fullwidth,language=python,numbers=left]
import pytest
from math_utils import divide

@pytest.fixture
def numerator():
    return 10.0

@pytest.fixture
def denominator():
    return 2.0

def test_divide(numerator, denominator):
    assert divide(numerator, denominator) == 5.0

def test_divide_type(numerator, denominator):
    result = divide(numerator, denominator)
    assert isinstance(result, float)
\end{lstlisting}
\end{column}
\begin{column}[c]{0.49\linewidth}
\begin{block}{}
\begin{itemize}
\item Decorate a function with \texttt{@pytest.fixture}.
\item Add the fixture name as a test parameter.
\item PyTest injects the return value automatically.
\item Multiple tests can share the same fixture.
\end{itemize}
\end{block}
\end{column}
\end{columns}
\end{frame}

\begin{frame}[fragile,squeeze]{Fixtures with Setup and Teardown}
\setlength{\leftmargini}{0.5cm}
\begin{columns}[c,t,totalwidth=1.0\textwidth]
\begin{column}[c]{0.49\linewidth}
\begin{lstlisting}[style=fullwidth,language=python,numbers=left]
import pytest, tempfile, os

@pytest.fixture
def temp_file():
    # Setup: runs before the test
    fd, path = tempfile.mkstemp()
    os.close(fd)

    yield path        # value provided to test

    # Teardown: always runs after the test
    if os.path.exists(path):
        os.remove(path)

def test_file_is_created(temp_file):
    assert os.path.exists(temp_file)

def test_file_is_writable(temp_file):
    with open(temp_file, 'w') as f:
        f.write("hello")
    assert os.path.getsize(temp_file) > 0
\end{lstlisting}
\end{column}
\begin{column}[c]{0.49\linewidth}
\begin{block}{}
\begin{itemize}
\item Use \texttt{yield} instead of \texttt{return} to add teardown.
\item Code \textit{before} \texttt{yield} = setup.
\item Code \textit{after} \texttt{yield} = teardown (always runs, even if the test fails).
\item Ideal for resources that need cleanup, e.g. files.
\end{itemize}
\end{block}
\end{column}
\end{columns}
\end{frame}


\section{Misc Features}

\begin{frame}[fragile,squeeze]{Testing Many Cases: \texttt{@pytest.mark.parametrize}}
\setlength{\leftmargini}{0.5cm}
\begin{columns}[c,t,totalwidth=1.0\textwidth]
\begin{column}[c]{0.49\linewidth}
\begin{lstlisting}[style=fullwidth,language=python,numbers=left]
import pytest
from math_utils import add

@pytest.mark.parametrize("a, b, expected", [
    (1,    2,    3),
    (0,    0,    0),
    (-1,   1,    0),
    (-2,  -3,   -5),
    (100, 200,  300),
])
def test_add(a, b, expected):
    assert add(a, b) == expected
\end{lstlisting}
Running produces:
\begin{lstlisting}[style=fullwidth,language=bash,numbers=none]
test_add[1-2-3]        PASSED
test_add[0-0-0]        PASSED
test_add[-1-1-0]       PASSED
test_add[-2--3--5]     PASSED
test_add[100-200-300]  PASSED
\end{lstlisting}
\end{column}
\begin{column}[c]{0.49\linewidth}
\begin{block}{}
\begin{itemize}
\item Runs the test once per parameter set.
\item First argument: comma-separated parameter names (as a string).
\item Second argument: list of tuples, one per test case.
\item PyTest generates a unique ID for each case.
\item Equivalent to 5 separate test functions!
\item Failures pinpoint exactly which inputs caused the problem.
\end{itemize}
\end{block}
\end{column}
\end{columns}
\end{frame}

\begin{frame}[fragile,squeeze]{\texttt{confttest.py} File}
\setlength{\leftmargini}{0.25cm}
\begin{columns}[c,t,totalwidth=1.0\textwidth]
\begin{column}[c]{0.45\linewidth}
\BulletList{
	\item \textit{conftest.py} is automatically loaded by PyTest.
	\item Used to store common functionality needed throughout tests.
	\item Advanced features: fixture scope, pytest plugins, hooks
}
\end{column}
\begin{column}[c]{0.45\linewidth}
\begin{lstlisting}[style=fullwidth,language=python,numbers=left]
import pytest

@pytest.fixture(scope="function")
def sample_data():
    """Fixture that provides sample data for tests."""
    return {"numbers": [1, 2, 3], "text": "pytest"}
\end{lstlisting}
\end{column}
\end{columns}
\end{frame}



%%%%%%%%%%%%%%%%%%%%%%%%%%%%%%%%%%%%%%%%%%%%%%%%%%%
\section{Summary}
%%%%%%%%%%%%%%%%%%%%%%%%%%%%%%%%%%%%%%%%%%%%%%%%%%%

\TwoColumnSlide{What We Covered}
{\BulletList{
\item \textbf{Why test}: catch bugs early, prevent regressions, document expected behavior.
\item \textbf{Installation}: \inlinecode{pip install pytest} or \inlinecode{conda install pytest}.
\item \textbf{Discovery}: files named \inlinecode{test_*.py}, functions named \inlinecode{test_*}.
\item \textbf{Running}: \inlinecode{pytest -v}, \inlinecode{-k}, \inlinecode{-x}.
\item \textbf{Assertions}: use \inlinecode{assert}; use \inlinecode{pytest.approx} for floats.
}}
{\AddPicture{actual_vs_unit.png}}
{}


\TwoColumnSlide{Acknowledgements}
{\BulletList{
	\item Claude for a first draft.
	\item PyTest documentation (\texttt{docs.pytest.org}).
	\item NSF for funding the REU.
	\item ISU and Ames National Lab
}}
{
	\AddPicture{nsf.jpg}
}
{}

\end{document}
